\section{Agent Runtime 与 Agent Loop}

\subsection{Agent Runtime：Agent 的运行时环境}

Agent Runtime 是 Agent 执行的"舞台"，它提供 Agent 运行所需的一切环境支持：

\begin{itemize}
    \item \textbf{Workspace}：如 OpenAI Codex 的工作空间
    \item \textbf{启动文件注入}：系统级的提示词（System Prompt）
    \item \textbf{内置 Tools 和 Skills}：预定义的工具集和技能文档
    \item \textbf{配置文件}：agents.md、src.md、tools.md 等 Markdown 配置文件
    \item \textbf{Memory 系统}：记忆和持久化机制
\end{itemize}

\subsection{Agent Loop：消息驱动的执行循环}

Agent Loop 描述了当一条消息真正进入系统后，如何被处理并最终生成响应或动作：

\begin{importantbox}{Agent Loop 的标准流程}
\begin{enumerate}
    \item \textbf{Intake}：接收用户消息（可能需要排队）
    \item \textbf{Prompt Injection}：大量的提示词注入（系统指令、Skills、环境信息等）
    \item \textbf{Model Inference}：调用远端或本地大模型进行推理
    \item \textbf{Function Call}（可选）：如果模型决定调用工具，执行对应的 Function
    \item \textbf{Streaming Response}：流式输出回复
    \item \textbf{Persistence}：持久化到本地（保存对话历史）
\end{enumerate}
这是以 OpenAI Codex 官方文档描述的 Agent Loop 为基础的标准流程。
\end{importantbox}

\begin{knowledgebox}{Agent Runtime vs. Agent Loop}
\begin{itemize}
    \item \textbf{Agent Runtime}：静态的运行时环境，在 Agent 启动时准备好，提供执行所需的一切资源
    \item \textbf{Agent Loop}：动态的执行循环，当消息进入后驱动 Agent 执行动作和生成回复
\end{itemize}
Runtime 是舞台，Loop 是舞台上演出的过程。
\end{knowledgebox}

\subsection{本章小结}

Agent Runtime 提供运行时环境（workspace、配置文件、内置工具），Agent Loop 是消息驱动的执行循环（接收→注入→推理→调用→响应→持久化）。两者共同构成 Modern Agent 的完整架构。


